\section*{第 \ref{ch:g12:metamorphic-rocks} 章——变质岩}

\begin{solution}{exo:g12:metamorphic-rocks:1}
(a) \omterm{def:g12:metamorphic-rocks:polymorph}{同质多象变体}；(b) \omterm{def:g12:metamorphic-rocks:index-mineral}{等变质线}；(c) \omterm{def:g12:metamorphic-rocks:contact-regional}{接触变质晕}；(d) \omterm{def:g12:metamorphic-rocks:foliation}{面理}；
(e) \omterm{def:g12:metamorphic-rocks:lithostatic}{静岩压力}。
\end{solution}

\begin{solution}{exo:g12:metamorphic-rocks:2}
(1) 大理岩；(2) 石英岩；(3) 板岩；(4) 片麻岩；(5) 角岩。
\end{solution}

\begin{solution}{exo:g12:metamorphic-rocks:3}
$P = \rho g z$：在 \qty{10}{km}，$2830 \times 9.81 \times \num{10000} \approx
\qty{2.78e8}{Pa} = \qty{0.278}{GPa}$；在 \qty{20}{km} 处则为两倍，
\qty{0.555}{GPa}。
\end{solution}

\begin{solution}{exo:g12:metamorphic-rocks:4}
(1) 蓝晶石（蓝晶石–红柱石分界线约在 \qty{0.25}{GPa}
、\qty{400}{\celsius} 处）。(2) 红柱石。(3) 夕线石。
\end{solution}

\begin{solution}{exo:g12:metamorphic-rocks:5}
(1) 板岩、千枚岩、片岩、片麻岩。(2) 在片岩中。(3) 片麻岩。
\end{solution}

\begin{solution}{exo:g12:metamorphic-rocks:6}
(1) $2830 \times 9.81 \times \num{35000} \approx \qty{9.72e8}{Pa}$，即
$\num{9.72e8}/101325 \approx \num{9600}$ 个大气压。(2) $1025 \times 9.81
\times 4000 \approx \qty{4.02e7}{Pa} = \qty{40.2}{MPa}$。(3) $972/40.2
\approx 24$ 倍。
\end{solution}

\begin{solution}{exo:g12:metamorphic-rocks:7}
(1) $z = \num{0.375e9}/(2830 \times 9.81) \approx \qty{13.5}{km}$。
(2) $(504 - 10)/13.5 \approx \qty{36.6}{\celsius}$ 每千米。(3) 否：
岩石须恰好处于图中的某一点；
这类岩石很少见。
\end{solution}

\begin{solution}{exo:g12:metamorphic-rocks:8}
(1) $23/3 \approx \qty{8}{mm}$。(2) 约 800 万年。(3) 岩石
在深部、高温条件下停留了数百万年。
\end{solution}

\begin{solution}{exo:g12:metamorphic-rocks:9}
(1) 四条\omterm{def:g12:metamorphic-rocks:index-mineral}{等变质线}，依次为石榴石、十字石、蓝晶石和夕线石，
(2) 向西。(3) 位于西端，处于夕线石带。其
岩石越过了蓝晶石–夕线石分界线：埋深约 \qty{20}{km}
（\qty{0.56}{GPa}），温度高于约 \qty{590}{\celsius}。
\end{solution}

\begin{solution}{exo:g12:metamorphic-rocks:10}
(1) 泥变为角岩，石灰岩变为大理岩，砂岩变为
石英岩。(2) 它们是在未受挤压的条件下被加热的：没有定向压力
使其矿物定向排列。(3) 其周围的岩石已经很热，因此
岩浆再使其升温的幅度很小。
\end{solution}

\begin{solution}{exo:g12:metamorphic-rocks:11}
(1) 垂直于挤压方向：其扁平面朝东西向，因此片状矿物
立于南北向的平面内。(2) 竖直，走向南北。
(3) 方解石和石英颗粒既不扁平也不细长，因此即便它们在挤压下重结晶，
也无从定向排列。
\end{solution}

\begin{solution}{exo:g12:metamorphic-rocks:12}
(1) $z = (T - 10)/30$，所以 $P = 0.0278\,(T - 10)/30 \approx
0.000927\,(T - 10)$。(2) $0.000927\,(T - 10) = 0.577 + 0.00209\,(T - 600)$
得 $0.001163\,T \approx 0.668$，故 $T \approx \qty{574}{\celsius}$，$P
\approx 0.000927 \times 564 \approx \qty{0.523}{GPa}$，且 $z = (574 - 10)/30
\approx \qty{18.8}{km}$。(3) 蓝晶石位于该深度以上，夕线石位于其以下。
\end{solution}

\begin{solution}{exo:g12:metamorphic-rocks:13}
(1) $10 + 10 \times 35 = \qty{360}{\celsius}$。(2) $10 + 30 \times 35 =
\qty{1060}{\celsius}$。(3) 在 \qty{35}{km} 处的压强约为
\qty{0.97}{GPa}；在 \qty{360}{\celsius} 时，该点深居蓝晶石场之中。
第二个点位于图的右侧远处，
处于夕线石场内，温度已达地壳岩石开始熔融的程度。
(4) 它们需要高压、低温的条件，这种组合只有
\omterm{def:g11:subduction-zones:subduction}{俯冲带}中被带入深部的冷地壳才能提供。
\end{solution}

\begin{solution}{exo:g12:metamorphic-rocks:14}
(1) 蓝晶石 $1000/3.61 \approx \qty{277}{cm^3}$；夕线石 $1000/3.24
\approx \qty{309}{cm^3}$；红柱石 $1000/3.15 \approx \qty{317}{cm^3}$。
(2) $(317 - 277)/317 \approx 0.127$，约 \qty{13}{\percent}。(3) 高压下
最紧密的原子排列更为有利；蓝晶石
以最小的空间容纳相同的原子。
\end{solution}

\begin{solution}{exo:g12:metamorphic-rocks:15}
(1) 在 \qty{20}{km} 处为 \qty{0.555}{GPa}，在 \qty{25}{km} 处为 \qty{0.694}{GPa}。
(2) $T = 600 + (P - 0.577)/0.00209$：约 \qty{590}{\celsius}（
\qty{0.555}{GPa}），\qty{656}{\celsius}（\qty{0.694}{GPa}）。
\qty{700}{\celsius} 时，岩石温度高于上述两者，处于夕线石场内：
结果一致。(3) $(700 - 10)/20 \approx \qty{34.5}{\celsius}$ 与
$(700 - 10)/25 \approx \qty{27.6}{\celsius}$ 每千米，接近典型的\omterm{def:g11:lithosphere-asthenosphere:geotherm}{地温梯度}。
(4) 由 $\qty{20}{km}$ 至 $\qty{25}{km}$，历时 4000 万年：
即每百万年 0.50 至 \qty{0.63}{km}，亦即每年 0.50 至
\qty{0.63}{mm}。
\end{solution}

\begin{solution}{pb:g12:metamorphic-rocks:1}
\textbf{1.} 该岩石具有\omterm{def:g12:metamorphic-rocks:foliation}{面理}，其中的云母大到可以肉眼看见。
白云母与黑云母构成了\omterm{def:g12:metamorphic-rocks:foliation}{片理}。

\textbf{2.} \ce{Al2SiO5}；它们互为\omterm{def:g12:metamorphic-rocks:polymorph}{同质多象变体}。

\textbf{3.} 每千克的体积之比为 $3.15/3.61 \approx
0.873$：蓝晶石约小 \qty{13}{\percent}。压强
有利于生成蓝晶石。

\textbf{4.} 蓝晶石带。岩石下沉过程中依次穿过了
黑云母、石榴石、十字石和蓝晶石的\omterm{def:g12:metamorphic-rocks:index-mineral}{等变质线}。

\textbf{5.} 约 800 万年。

\textbf{6.} $P = 0.577 + 0.00209 \times 0 = \qty{0.577}{GPa}$.

\textbf{7.} \qty{0.70}{GPa} 高于 \qty{0.577}{GPa}，位于
蓝晶石–夕线石分界线的蓝晶石一侧；它也远高于红柱石区，
该区在 \qty{600}{\celsius} 时位于约 \qty{0.23}{GPa} 以下。

\textbf{8.} \qty{0.577}{GPa}.

\textbf{9.} 由 $0.70 = 0.577 + 0.00209\,(T - 600)$ 可得 $T = 600 +
0.123/0.00209 \approx \qty{659}{\celsius}$。

\textbf{10.} \qty{7.0}{kbar}；$\num{7.0e8}/101325 \approx \num{6900}$
个大气压。

\textbf{11.} $\rho g = 2830 \times 9.81 \approx \qty{27800}{Pa}$ 每米，
即 \qty{27.8}{MPa} 每千米。

\textbf{12.} $z = \num{0.70e9}/\num{27762} \approx \qty{25200}{m}$，约
\qty{25.2}{km}。

\textbf{13.} $\num{0.577e9}/\num{27762} \approx \qty{20.8}{km}$.

\textbf{14.} $\num{0.70e9}/(2700 \times 9.81) \approx \qty{26.4}{km}$，约
深 \qty{5}{\percent}（$26.4/25.2 \approx 1.048$）。

\textbf{15.} 位于地壳内：\qty{25}{km} 甚至小于正常地壳
厚度。

\textbf{16.} $(600 - 10)/25.2 \approx \qty{23.4}{\celsius}$ 每千米。

\textbf{17.} 比\omterm{def:g11:subduction-zones:subduction}{俯冲带}（每千米低于 \qty{10}{\celsius}）高得多，
比典型地壳（\qty{30}{\celsius}）略低，
远离火山区（40 至 \qty{50}{\celsius}）：一块
埋藏于山脉加厚地壳深处的岩石，
经\omterm{def:g12:metamorphic-rocks:contact-regional}{区域变质作用}形成。

\textbf{18.} $10 + 30 \times 25.2 \approx \qty{766}{\celsius}$。蓝晶石--
夕线石分界线延长后，在该温度下位于 $0.577 + 0.00209 \times 166
\approx \qty{0.92}{GPa}$ 处；\qty{0.70}{GPa} 低于该值，
位于夕线石稳定区内。若按典型\omterm{def:g11:lithosphere-asthenosphere:geotherm}{地温线}，本应形成夕线石。

\textbf{19.} $\qty{25.2}{km} = \qty{2.52e7}{mm}$；按每年 \qty{0.5}{mm} 计，
约 5000 万年。

\textbf{20.} 岩石上升时冷却，反应因此减慢，
而逆向反应所需的水在下沉途中
已被驱离。

\textbf{21.} 蓝晶石片岩结晶于约 \textbf{\qty{25}{km}
深处}，温度约 \qty{600}{\celsius}。
\end{solution}
